\section*{INTRODUCCIÓN}
\addcontentsline{toc}{section}{INTRODUCCIÓN}

La transición global hacia matrices energéticas descarbonizadas y sostenibles ha dado inicio a una transformación estructural en la arquitectura de los sistemas eléctricos de potencia: a finales de 2025 las fuentes renovables representaban el $49$ por ciento de la capacidad instalada mundial y el $85.6$ por ciento de las adiciones anuales de capacidad \citep{IRENA2026}. En este nuevo paradigma, las microrredes eléctricas emergen como subsistemas estratégicos, capaces de integrar de manera flexible y coordinada Fuentes No Convencionales de Energía Renovable (FNCER), como la generación solar fotovoltaica, los parques eólicos y los sistemas de almacenamiento de energía en baterías (BESS), operando tanto en modo interconectado a la red de distribución como en modo isla o aislado \citep{Mamun2026, Altin2023MultiAgent}.

No obstante, la sustitución progresiva de generadores síncronos convencionales por recursos energéticos distribuidos basados exclusivamente en convertidores electrónicos de potencia (IBRs, \textit{Inverter-Based Resources}) impone retos operacionales críticos sobre la estabilidad electromecánica. Al carecer de masas giratorias acopladas magnéticamente, los convertidores estáticos reducen la inercia rotacional equivalente del sistema ($H \to 0$), provocando incrementos en la tasa de cambio de la frecuencia (\textit{Rate of Change of Frequency}, RoCoF) ante el menor desbalance entre generación y demanda \citep{Mehrzad2025}. Simultáneamente, en las redes de distribución de media y baja tensión, la elevada relación resistiva-inductiva de los alimentadores ($R/X \ge 1$) genera un fuerte acoplamiento cruzado no lineal entre las potencias activa y reactiva ($P-Q$) con la magnitud y fase de la tensión ($V-\theta$), lo cual agrava la aparición de desviaciones nodales \citep{Sowa2023, Pompodakis2022Sensitivity}.

Para gestionar estas fluctuaciones en la escala de control primario, los inversores implementan esquemas descentralizados de control por estatismo (\textit{droop control}); sin embargo, esta estrategia introduce desviaciones inevitables en estado estacionario tanto en la frecuencia como en el perfil de voltaje nominal \citep{Jasim2023}. Aunque las arquitecturas de control secundario centralizado han sido la solución convencional, su implementación en microrredes modernas resulta inviable debido a su alta vulnerabilidad frente a puntos únicos de falla (\textit{Single Point of Failure}, SPoF), la escasa escalabilidad y los cuellos de botella telemétricos en los canales de comunicación \citep{Khayat2020, Altin2023MultiAgent}.

Frente a este desafío, la adopción de Sistemas Multiagente (MAS) operando bajo protocolos de consenso distribuido ha adquirido un notable protagonismo en la literatura científica \citep{Mamun2026, Altin2023MultiAgent}. En un esquema MAS, la toma de decisiones se descentraliza hacia la periferia de la red (\textit{Edge Computing}), donde cada recurso distribuido es gobernado de forma autónoma por un agente inteligente que procesa información local y se coordina exclusivamente con sus nodos vecinos directos \citep{Diaz2025, AlvarezLopez2026}. Sin embargo, la revisión del estado del arte revela una brecha metodológica: la polarización entre simuladores de potencia puros (que modelan la física pero ignoran la asincronía y protocolos de red) y los marcos MAS abstractos (que modelan la cognición y la topología de comunicación pero tienden a resolver el circuito con balances algebraicos, sin anclaje físico o \textit{physical grounding}) \citep{Muszynski2026EnergyTwin, Khanna2025}.

El presente trabajo de grado aborda esta problemática mediante el diseño, implementación y validación de una arquitectura de control secundario distribuido basada en \textbf{consenso en tiempo fijo (FxTS)} sobre topologías de grafo, cuya ley incorpora términos correctores no lineales que garantizan una cota de establecimiento independiente de las condiciones iniciales, y soportada por una plataforma de \textbf{Co-Simulación Ciber-Física Multitasa} \citep{Gomes2022FMI30, Liu2025Multirate}, desplegada en hardware de borde real —un clúster de dispositivos Raspberry Pi 5 con contenedores Docker modulares, en línea con el paradigma de computación en el borde para redes inteligentes \citep{Mehmood2021}—. La plataforma desacopla la resolución de los transitorios rápidos de la electrónica de potencia ($h = 1.0\text{ ms}$) del solucionador macroscópico de flujo de potencia de la red ($H = 500\text{ ms}$), integrando sockets síncronos ZeroMQ para la sincronización temporal y el consenso P2P, además de la implementación del protocolo telemétrico asíncrono MQTT para la inyección de contingencias climáticas y perturbaciones de red \citep{Khanna2025, Haid2022CosimCorrector, Kyllingstad2025CosimErrors}.

El documento de tesis se encuentra estructurado en ocho secciones fundamentales:
\begin{itemize}
    \item \textbf{Planteamiento del Problema:} Expone la crisis de inercia en microrredes aisladas, el conflicto metodológico del anclaje físico y formula la pregunta de investigación central.
    \item \textbf{Objetivos:} Define el objetivo general y los cuatro objetivos específicos orientados a la concepción matemática, diseño de co-simulación, implementación en hardware embebido y evaluación cuantitativa de resiliencia.
    \item \textbf{Justificación y Alcance:} Sustenta la pertinencia teórica, metodológica, tecnológica y socioeconómica del proyecto, delimitando sus fronteras experimentales.
    \item \textbf{Antecedentes y Estado del Arte:} Analiza críticamente los avances en control secundario por consenso, inteligencia en el borde, resiliencia telemétrica y co-simulación multitasa.
    \item \textbf{Marco Teórico:} Desarrolla exhaustivamente las ecuaciones dinámicas de los convertidores $dq$, la sintonización del SRF-PLL, los solucionadores de flujo de potencia (FBS y Matrices de Sensibilidad), y la demostración de estabilidad de Lyapunov para el consenso en tiempo fijo.
    \item \textbf{Diseño Metodológico:} Detalla las cuatro fases de ingeniería y los resultados esperados del proyecto.
    \item \textbf{Arquitectura e Implementación del Sistema Multiagente y Co-Simulación:} Expone la plataforma ciber-física, el motor de co-simulación multitasa, los modelos dinámicos de los cinco recursos distribuidos y el despliegue en hardware de borde.
    \item \textbf{Resultados y Discusión:} Presenta la validación en dos niveles —benchmark canónico multicriterio de control secundario y banco físico HIL sobre clúster de Raspberry Pi 5—, el contraste de los resultados con las cotas teóricas del marco teórico y la discusión integrada de los límites de capacidad y de las amenazas a la validez.
\end{itemize}